\subsection{The second marked elliptic cover and the exact joint arithmetic fibre}

\begin{statement}{D.36}{Proved}{A quartic cover, a shared degree-twelve cover, and the retained CM discrepancy}
This addresses Part D(e) and the two-torus follow-up to LM2,
manuscript \S51.15, folios 725--727, scans 156--160, and
\S51.16, folio 734 verso, scan 174. Write
$U=\mathbb P^1\setminus\{0,1,\infty\}$,
$F=\pi_1(U)=\langle x,y\rangle$, and let $K_n$ be the kernel of
$F\to\mathbb Z/n$ sending both $x,y$ to one. Its marked cover is
\[
 C_n^\circ:\quad v^n=z(z-1),\qquad z\in U.             \tag{D.36.1}
\]
Here $C_3=\mathcal E_0$ of D.34. The smooth compactification of
$C_4^\circ$ is an elliptic curve with four removed points. The
following rational maps are inverse isomorphisms of compact curves:
\[
 \begin{gathered}
 \psi_0:C_4\longrightarrow E_4,\qquad E_4:Y^2=X^3-X,\\
 (z,v)\longmapsto
 \left({z\over v^2},{v\over z-1}\right),\qquad
 (X,Y)\longmapsto
 \left({X^2\over X^2-1},{Y\over X^2-1}\right).
 \end{gathered}                                      \tag{D.36.2}
\]
The four omitted points map exactly to $E_4[2]$.
Translation by $(0,0)$ after $\psi_0$ sends the point over $z=0$
to the origin and is defined over $\mathbb Q$.
The deck map $v\mapsto iv$ becomes $[i]:(X,Y)\mapsto(-X,iY)$;
the swap $z\mapsto1-z$ becomes $T_{(0,0)}[-1]$.

Restriction to $\widehat K_n$ and filling all cusp subgroups give
actual maps from automorphisms of $\widehat F$ preserving each of
the three ordered cusp conjugacy classes with one common multiplier
to the homology of $C_n$. For $n=3,4$ their compact targets are precisely
$T\mathcal E_0$ and $TE_4$ with the marked isomorphisms above.
With the positive $z$ tangent and the lift
$v=\zeta_8 z^{1/4}(1-z)^{1/4}$, the restricted Galois action on
the quartic compact Tate module is
\[
 \rho^{\rm res}_4(g)=[i]^{\kappa_{-1,4}(g)}\rho_{E_4}(g),
 \qquad\kappa_{-1,4}(g)={\chi_8(g)-1\over2}\pmod4.
                                                               \tag{D.36.3}
\]
The cubic lift with leading coefficient $-1$ is rational and has no
corresponding tangent correction. Over
$L=\mathbb Q(\zeta_{12})$, the factor in (D.36.3) is a global sign
and disappears on projective three-torsion.

The common cyclic cover and its two maps are
\[
 C_{12}\longrightarrow C_3,\ (z,v)\mapsto(z,v^4),
 \qquad
 C_{12}\longrightarrow C_4,\ (z,v)\mapsto(z,v^3).
                                                               \tag{D.36.4}
\]
It has genus five and four punctures. On compact integral homology
the joint map is onto and has kernel $\mathbb Z^6$:
\[
 0\longrightarrow\mathbb Z^6\longrightarrow
 H_1(C_{12},\mathbb Z)
 \longrightarrow H_1(C_3,\mathbb Z)\oplus H_1(C_4,\mathbb Z)
 \longrightarrow0.                                   \tag{D.36.5}
\]
All four nonzero Smith invariants are one. The $S_3$ normal core of
the cyclic-twelve character has deck group
$(\mathbb Z/4)^2\times\mathbb Z/3$, of order $48$; the full
Fermat-twelve abelian cover has deck group $(\mathbb Z/12)^2$,
of order $144$, and maps to this core with kernel of order three.

There is an exact arithmetic joint image over $L$. Put
$\mathcal O=\mathbb Z_2[\zeta_3]$. For a norm-one-modulo-four
unit $\eta\in(\mathcal O/4)^\times$, let $\tau(\eta)\in\mu_3$
be the unique element with the same reduction modulo two, and set
\[
 \epsilon(\eta)=\tau(\eta)^{-1}\eta\in\{I,-I\}.
                                                               \tag{D.36.6}
\]
Then the actual image of
$(\rho_{\mathcal E_0,4},\mathbb P\rho_{E_4,3})(G_L)$ is exactly
the six-element graph
\[
 \left\{\left(\eta,
 [i]^{(1-\epsilon(\eta))/2}\right):
 N(\eta)=1\pmod4\right\}.                            \tag{D.36.7}
\]
The exponent treats $\epsilon$ as the scalar sign $\pm1$.
The joint fixed field is
$L(\sqrt[3]{4},r_0)$, of degree six over $L$, where
$r_0=(3/4)^{1/4}$; the common subfield of the two torsion fields
is $L(r_0)$, of degree two. The normalized unit (D.36.6) is precisely the
actual $E_3[4]$ action after D.34's cubic twist is removed.

For the family retained in D.33 these three specified coordinates are
\[
 \epsilon_{\rm cover}=-1,\qquad
 T_4^{\rm LM2}=I,\qquad W_3^{\rm LM2}=[i].             \tag{D.36.8}
\]
The first sign is unchanged by every Ihara gauge
$\operatorname{Int}(x^k)$: this multiplies its cover Tate action by
$\zeta_3^k$, while the exact simultaneous change
$(A,\beta,V)\mapsto(x^kA,x^k\beta,x^kV)$ fixes
$T=A^{-1}V$ and $W=\beta^{-1}V$.
The actual norm-one Galois element supplied by D.35 can realize the
entire D.33 compact Tate-sixteen matrix and the displayed $W_3$;
its actual $T_4$ is $-I$. Thus the full prescribed finite action is
still an additional requirement. Equality of these cover and
fixed-fibre signs for every genuine full profinite
$\widehat{\rm GT}_0$ element has not been proved here.
\end{statement}
\begin{proof}
Substitution of (D.36.2) gives
$Y^2-X^3+X=0$, using $v^4=z(z-1)$, and substitution of its
inverse gives the original equation. Both rational maps are inverse
on dense open sets and therefore extend to inverse maps of the
smooth projective curves. At $z=0$, $v$ is a local parameter,
$z=-v^4+O(v^8)$, and the image is $(0,0)$. At $z=1$ the
coordinates have pole orders two and three, giving the origin.
The two points over infinity have $z/v^2\to\pm1$ and
$v/(z-1)\to0$, giving $(1,0),(-1,0)$. These are exactly the four
two-torsion points. The point $(0,0)$ has order two, so the stated
translation gives the pointed isomorphism over $\mathbb Q$.

The two deck calculations are exact substitutions. Multiplying $v$
by $i$ sends $(X,Y)$ to $(-X,iY)$. Replacing $z$ by $1-z$ sends
it to $(-1/X,-Y/X^2)$. The addition law on $E_4$ gives
$T_{(0,0)}(X,Y)=(-1/X,Y/X^2)$, proving the swap formula. Since
$[i]$ fixes $(0,0)$, these formulas also hold after the pointed
translation.

We describe the source and all integral quotient maps explicitly.
Choose coset representatives $1,x,\ldots,x^{n-1}$ for $K_n$.
Its Schreier basis consists of
\[
 a=x^n,\quad b_j=x^j yx^{-(j+1\bmod n)},
                    \qquad0\leq j<n.                 \tag{D.36.9}
\]
The peripheral words are $x^n$, $y^n$, and
$x^j((xy)^{-1})^{n/d}x^{-j}$ for $0\leq j<d$, where
$d=\gcd(n,2)$. Their $d+2$ homology vectors sum to zero and span
a primitive rank-$d+1$ lattice. Thus the compact quotient has rank
$n-d$ and genus $(n-d)/2$. This proves all stated genera,
puncture counts and free ranks. It also proves the exact profinite
filling map by killing the normal closures of these actual cusp
words. An automorphism preserving the three ordered cusp classes acts
on $F^{\rm ab}$ by the common scalar unit, so it preserves
$\widehat K_n$ and the family of lifted
cusp subgroups. Restriction and filling are consequently defined
for that entire pure source, without a Galois realization assumption.
The swap also preserves the cyclic character $(1,1)$; a three-cycle
does not preserve $K_4$, as it sends that character to $(1,-2)$.
The later normal-core construction is the precise enlargement on
which the complete $S_3$ action is defined.

The quartic arithmetic tangent factor follows within the arithmetic
fundamental group. Let $s_z$ be its section at the positive $z$
tangent. In the displayed Puiseux lift, $s_z(g)$ multiplies the
leading coefficient by
$g(\zeta_8)/\zeta_8=i^{\kappa_{-1,4}(g)}$.
The positive loop $x$ multiplies $v$ by $i$. Therefore the section
at the rational $v$ tangent on the cover is
$x^{-\kappa_{-1,4}(g)}s_z(g)$ modulo $\widehat K_4$.
Changing the chosen lift of this exponent only changes the section
by a member of $\widehat K_4$, whose inner action is trivial on
compact homology. Conjugation and the deck formula $x\mapsto[i]$
give (D.36.3), with the displayed sign of the exponent. For the
cubic cover the lift is
$v=-z^{1/3}(1-z)^{1/3}$, whose leading coefficient is rational.
For $g\in G_L$, $\chi_8(g)$ is $1$ or $5$ modulo eight, so the
quartic factor is $I$ or $-I$. Its projective action is therefore
the actual projective $E_4$ action.

Inclusion of the actual Schreier words (D.36.9) for $n=12$ in
those for $n=3,4$ gives the homology maps in (D.36.4). The
certificate records their punctured matrices and primitive cusp
quotients. In its compact bases the two matrices are
\[
 \begin{gathered}
 M_3=\begin{pmatrix}
 1&0&-1&1&0&-1&1&0&-1&1\\
 0&1&-1&0&1&-1&0&1&-1&0
 \end{pmatrix},\\[4pt]
 M_4=\begin{pmatrix}
 1&0&-1&0&1&0&-1&0&1&0\\
 0&1&0&-1&0&1&0&-1&0&1
 \end{pmatrix}.
 \end{gathered}                                      \tag{D.36.10}
\]
The common cover uses compatible lifts $v^4,v^3$; converting its
cubic leading coefficient $\zeta_6$ to D.34's coefficient $-1$
is exactly the deck map $\delta$. This is defined over $L$ and
does not alter the scalar sign (D.36.6).
The stack of the two matrices has rank four and Smith diagonal
$(1,1,1,1)$, proving (D.36.5). Its deck-equivariance is checked
against the full order-twelve matrix and the elliptic matrices
\[
 D_3=\begin{pmatrix}0&-1\\1&-1\end{pmatrix},\qquad
 D_4=\begin{pmatrix}0&-1\\1&0\end{pmatrix}.
                                                               \tag{D.36.11}
\]
The $S_3$ orbit of the cyclic character is generated, up to units,
by $(1,1),(1,-2),(-2,1)$. Modulo four their common kernel is zero.
Modulo twelve it is exactly
$\{(0,0),(4,8),(8,4)\}$. This proves the normal-core orders and
the actual quotient from the Fermat-twelve deck group. It does not
replace the smaller normal core by the full abelian cover.

For the arithmetic graph, D.35 proves
$\rho_{\mathcal E_0}(G_{\mathbb Q(\zeta_3)})=\mathcal O^\times$.
The subgroup fixing $L=\mathbb Q(\zeta_3,i)$ is cut out by
$\chi_4=1$, equivalently by the norm modulo four. Hence its image
on $\mathcal E_0[4]$ is exactly the six-element norm kernel in
$(\mathcal O/4)^\times$. Reduction modulo two identifies its
unique order-three factor. D.34's formula
$\rho_{\mathcal E_0}=\zeta_3^{-\kappa_{4,3}}\rho_{E_3}$ and
the trivial $G_{\mathbb Q(\zeta_3)}$ action on $E_3[2]$ show that
this factor is exactly $\tau=\zeta_3^{-\kappa_{4,3}}$.
The residual principal unit has norm one modulo four and is
therefore exactly $I$ or $-I$. This proves (D.36.6) as an equality
with the actual $E_3[4]$ representation, including its cubic twist.

D.25 gives the common field $M'=\mathbb Q(i,r_0)$ for full
$E_3[4]$ and projective $E_4[3]$. Over $L$ it is quadratic,
and its nontrivial automorphism sends $r_0$ to $-r_0$.
On $E_3[4]$ this is $-I$. To identify its other action without a
matrix-sign assumption, put $w_0=1/r_0$ and use the exact
three-division roots from D.25:
\[
 a={w_0+2/w_0\over2},\qquad
 b={w_0-2/w_0\over2i}.                                \tag{D.36.12}
\]
The involution sends $w_0$ to $-w_0$ and consequently sends every
$x$-coordinate $\pm a,\pm b$ to its negative. This is precisely
the projective action of $[i]$ on $E_4[3]$. For a nonzero point
$P\in E_4[3]$, the basis $(P,[i]P)$ realizes it by the matrix
$D_4\bmod3$; independence follows from the irreducibility of
$X^2+1$ over $\mathbb F_3$. This proves the graph (D.36.7).
The cubic extension $L(\sqrt[3]{4})/L$ is independent of the
quadratic extension $L(r_0)/L$, and the product map to the two
distinct factors $\mu_3$ and $\{\pm I\}$ is injective.
Its kernel field is therefore exactly $L(\sqrt[3]{4},r_0)$.

D.34 computes the D.33 compact matrix modulo sixteen; modulo four
it is $-I$, so its order-three factor is one and its normalized sign
is $-1$. D.25--D.33 retain $A_{12}=V_{12}=I$ and
$\beta_{12}=(\begin{smallmatrix}9&8\\4&9\end{smallmatrix})$
as the fixed CM coordinates, giving $T_4=I$ and
$W_3=[D_4]$. The gauge comparison is exact: $x=t^{-2}$ acts
on the cubic compact homology as $\delta=[\zeta_3]$, while its
simultaneous left multiplication on $A,\beta,V$ cancels in $T,W$.
Removing the unique order-three factor in (D.36.6) cancels that
same $\zeta_3^k$ and leaves the cover sign unchanged. This proves
every part of (D.36.8), including its stability under the specified
gauge. The norm-one unit of D.35 has this full Tate-sixteen matrix;
an actual Galois element realizing it fixes $L$ and consequently
has $W_3=[i]$ and $T_4=-I$ by (D.36.7).

We record the source checks that delimit the remaining transfer.
Nakamura, Tsunogai and Yasuda
\cite[Proposition 4.4]{dNTY2010} computes a beta-function cover
action for full $\widehat{\rm GT}$.
Their fixed-fibre comparison, Theorem 4.16, and braid
power-conjugacy Proposition 4.13 are stated for $G_{\mathbb Q}$;
the introduction, p.432, expressly leaves the broader implication
undecided. The later full mapping-class-group actions in
\cite[Theorem 1.1]{dLNS2026} come with an explicit qualification
on p.2 concerning compatibility with the gluing rules. These opened
results do not supply the missing equivariant comparison for the
specified $\widehat{\rm GT}_0$ source. The shared cover and the
exact finite coordinates above provide a concrete further test;
no nonempty full-source fibre follows from the arithmetic graph.
\end{proof}
